\documentclass[12pt]{article}
\usepackage{stfx_syllabus}

\begin{document}

\stfxsyllabusheader{Introduction to Classical Syriac}{RELS 4XX}{Winter 2027}

\section*{Welcome to Classical Syriac!}

\instructorinfo{Ken Penner}{Dr.\ Penner}{he/him/his}{kpenner@stfx.ca}{Nicholson Tower 503\newline Book via Outlook Bookings link on Kwe' (Moodle)}

\textbf{Contacting your instructor:} If you have any questions, you can contact me by email. Please include the course code (e.g., ``RELS 4XX'') and the topic in the subject line and sign the message with your full name and student number. I will try to respond to your email within 24 hours, but it could take 48 hours or longer for you to get a response on weekends and holidays, and in exceptional circumstances.

\section*{StFX University Land Acknowledgement}

StFX is located in Mi’kma’ki, the unceded ancestral territory of the Mi'kmaw people. The Mi'kmaw name for Antigonish is Nalikitquniejk, meaning ‘place where branches are torn off.’

Ms't wiaqpulti'kl ankukamkewe'l. We are all treaty people.

\section*{Course Overview}

Syriac, the Aramaic dialect of Edessa, was the language of one of the three great early Christian cultures, alongside Greek and Latin. Its literature includes the Peshitta Bible, Ephrem's hymns, Aphrahat's Demonstrations, the Acts of Thomas, and a vast body of saints' lives, homilies and chronicles. This course teaches you to read that literature.

The course follows one continuous story, \emph{The Acts of Mar Yukhannan}, written for this course in the style of the apocryphal Acts. Each of the thirteen lessons (0--12) adds grammar and about thirty high-frequency words, and then advances the story. By the end you will have read the whole tale in Syriac script and met every major verb stem.

No previous Semitic language is required, but students who know Hebrew, Aramaic or Arabic will recognize much. Pronunciation follows the West Syriac (Serto) tradition.

\subsection*{Course Outcomes}

By the end of this course, you should be able to:
\begin{itemize}
  \item read and write Syriac in the Serto script, with and without vowel signs;
  \item recognize and translate the nominal system: states, gender, number, pronominal suffixes and prepositions;
  \item parse verbs in the Peal, Paael and Aphel stems and their passives, including the common weak roots;
  \item use a lexicon (SEDRA or Payne Smith) to find words by root;
  \item translate simple narrative prose, with help, from unseen Syriac texts.
\end{itemize}

\subsection*{Course Materials}

\subsubsection*{Required Materials}

\emph{The Acts of Mar Yukhannan: A First Reader in Classical Syriac} (course textbook, PDF on Kwe'), with its paradigm tables (Appendix A) and glossary (Appendix B). Each lesson also has an interactive version on Kwe' with self-checking exercises.

\subsubsection*{Recommended}

J. F. Coakley, \emph{Robinson's Paradigms and Exercises in Syriac Grammar}, 6th ed. (Oxford, 2013); the SEDRA online lexicon (sedra.bethmardutho.org).

\section*{Important Information for Students}

\subsection*{Academic Integrity}

The highest standard of academic integrity is expected. All students must understand the meaning and consequences of academic offences such as plagiarism, cheating, tampering, and falsification. For more information, please see \href{https://www.stfx.ca/student-services/academic-services/academic-success-centre/academic-integrity}{Academic Integrity Policies and Procedures}.

\subsection*{Use of Artificial Intelligence (AI)}

For this course, the use of artificial intelligence (AI) aids is strictly prohibited for exams and graded forums unless explicitly authorized by the instructor.

\subsection*{Exam and Quiz Dates Policy}

Every exam and quiz must be written on the date, and within the time frame, stipulated in the Schedule posted on Kwe’ (Moodle) or the \href{https://www.stfx.ca/applications-admissions/registrars-office/examinations}{Exam Schedule set by the Registrar’s Office} (as applicable).

\subsection*{Withdrawal \& Refund Rules}

There are deadlines for withdrawing from this course, and possible academic consequences for doing so. See \href{https://drive.google.com/file/d/1PyNFH_oOmuxQtN2hc7oxZXAcQaHOnOAJ/view}{StFX Refund Policy} for more information.

\section*{Course Evaluation}

\begin{stfxevaluation}

\evalrow{Interactive Lessons}{
  \begin{evalitemlist}
    \item Completion of each lesson's self-checking exercises on Kwe' before class.
  \end{evalitemlist}
}{10}

\evalrow{Vocabulary Quizzes}{
  \begin{evalitemlist}
    \item Weekly short quizzes on the lesson's vocabulary; lowest two dropped.
  \end{evalitemlist}
}{10}

\evalrow{Homework and Translation}{
  \begin{evalitemlist}
    \item Written exercises and story translations, submitted weekly.
  \end{evalitemlist}
}{15}

\evalrow{Module Tests (2)}{
  \begin{evalitemlist}
    \item Module 1 test (Lessons 0--4) and Module 2 test (Lessons 5--8), 12.5\% each.
  \end{evalitemlist}
}{25}

\evalrow{Final Exam}{
  \begin{evalitemlist}
    \item Cumulative, including an unseen translation passage. Scheduled by the Registrar (April 9--21).
  \end{evalitemlist}
}{40}

\evaltotalrow{100}

\end{stfxevaluation}

\section*{Student Supports}

Your success as a student is supported by a comprehensive network of campus services:

\begin{itemize}
  \item \href{https://www.stfx.ca/student-services/academic-services/accessible-learning}{Tramble Centre for Accessible Learning}
  \item \href{https://www.stfx.ca/programs-courses/stfx-online/student-support}{Where to go for help?}
  \item \href{https://www.stfx.ca/student-life-support/student-services-department}{Student Services}
  \item \href{https://www.stfx.ca/applications-admissions/financial-support/financial-aid-office}{Financial Aid}
  \item \href{https://www.stfx.ca/student-services/support-services/visible-at-x}{Visible @ X}
  \item \href{https://www.stfx.ca/student-services/support-services/human-rights-equity}{Human Rights and Equity}
\end{itemize}

\section*{Course Policies}

\subsection*{Equitable Learning Environment}

Learning is more effective in a respectful, inclusive, and equitable environment. Instructors and students should collaborate to create learning spaces that uphold human dignity, equity, diversity, respect for the individual, and the open exchange of ideas.

\subsection*{Preferred Pronouns}

StFX respects the names and pronouns of all students, staff, and faculty. If you use a name that is different from your legal name, you can update your name by \href{https://stfx.teamdynamix.com/TDClient/1764/Portal/Requests/Service/39840/Change-of-Name}{submitting a request}.

\end{document}
